\section*{Capítulo \ref{ch:g12:stereochemistry} --- Estereoquímica: quiralidade e isômeros}

\begin{solution}{exo:g12:stereochemistry:1}
Butan-2-ol: o carbono 2 (\ce{H}, \ce{OH}, \ce{CH3}, \ce{CH2CH3}).
Propan-2-ol: nenhum (o carbono 2 carrega dois \ce{CH3} idênticos).
2-Clorobutano: o carbono 2. Ácido lático: o carbono 2 (\ce{H}, \ce{OH},
\ce{CH3}, \ce{COOH}).
\end{solution}

\begin{solution}{exo:g12:stereochemistry:2}
Primeiro: os dois \ce{Cl} (prioritários sobre o \ce{H}) do mesmo lado:
\omterm{def:g12:stereochemistry:z-e}{Z}. Segundo: à esquerda, o \ce{CH3} tem prioridade, em cima; à direita, o
\ce{CH2CH3} tem prioridade, embaixo: \omterm{def:g12:stereochemistry:z-e}{E}.
\end{solution}

\begin{solution}{exo:g12:stereochemistry:3}
\omterm{def:g12:stereochemistry:chiral}{Quirais}: a luva, o parafuso. Não \omterm{def:g12:stereochemistry:chiral}{quirais}: a colher, o cubo (cada um tem
um plano de simetria). O \ce{CHFClBr} é \omterm{def:g12:stereochemistry:chiral}{quiral} (quatro \omterm{def:g7:atoms-and-molecules:atom}{átomos} diferentes
no carbono); o \ce{CH2Cl2} não é.
\end{solution}

\begin{solution}{exo:g12:stereochemistry:4}
Uma \omterm{def:g6:pure-substances-and-mixtures:mixture}{mistura} de quantidades iguais de dois \omterm{def:g12:stereochemistry:enantiomer}{enantiômeros}. A carvona
racêmica contém os dois \omterm{def:g12:stereochemistry:enantiomer}{enantiômeros}, logo o nariz recebe os dois
cheiros ao mesmo tempo.
\end{solution}

\begin{solution}{exo:g12:stereochemistry:5}
Não: o carbono 1 carrega dois \omterm{def:g7:atoms-and-molecules:atom}{átomos} idênticos, \ce{H} e \ce{H};
trocá-los dá a mesma \omterm{def:g7:atoms-and-molecules:molecule}{molécula}.
\end{solution}

\begin{solution}{exo:g12:stereochemistry:6}
\setchemfig{atom sep=2em}
\chemfig{C(-[2]CH_2CH_3)(-[4]H_3C)(<[:-30]OH)(<:[:-150]H)} \quad espelho
\quad \chemfig{C(-[2]CH_2CH_3)(-[0]CH_3)(<[:-150]HO)(<:[:-30]H)}
\end{solution}

\begin{solution}{exo:g12:stereochemistry:7}
2-Clorobutano: 2 (um \omterm{def:g12:stereochemistry:chiral}{carbono assimétrico}). 2,3-Diclorobutano: 3 (dois
\omterm{def:g12:stereochemistry:chiral}{carbonos assimétricos} com metades iguais: um par de \omterm{def:g12:stereochemistry:enantiomer}{enantiômeros} e uma
forma meso). Pent-2-eno: 2 (\omterm{def:g12:stereochemistry:z-e}{Z} e \omterm{def:g12:stereochemistry:z-e}{E}).
\end{solution}

\begin{solution}{exo:g12:stereochemistry:8}
\omterm{def:g12:stereochemistry:enantiomer}{Enantiômeros}; \omterm{def:g12:stereochemistry:diastereomer}{diastereoisômeros}; \omterm{def:g12:stereochemistry:diastereomer}{diastereoisômeros}; a mesma \omterm{def:g7:atoms-and-molecules:molecule}{molécula} (o
ácido meso-tartárico é sobreponível à sua imagem no espelho).
\end{solution}

\begin{solution}{exo:g12:stereochemistry:9}
Na \omterm{def:g12:stereochemistry:conformation}{conformação} alternada com os grupos \ce{CH3} opostos, os grupos
grandes estão o mais longe possível um do outro: é a mais estável. Na
eclipsada, os dois \ce{CH3} ficam frente a frente e se atrapalham.
\end{solution}

\begin{solution}{exo:g12:stereochemistry:10}
As suas duas metades, cada uma \ce{CH(OH)COOH}, são imagens no espelho
uma da outra: um plano entre os dois carbonos centrais corta a \omterm{def:g7:atoms-and-molecules:molecule}{molécula}
em duas metades espelhadas. Uma \omterm{def:g7:atoms-and-molecules:molecule}{molécula} com um plano de simetria é a
sua própria imagem no espelho: não é \omterm{def:g12:stereochemistry:chiral}{quiral}.
\end{solution}

\begin{solution}{exo:g12:stereochemistry:11}
\setchemfig{atom sep=1.7em}
(\omterm{def:g12:stereochemistry:z-e}{Z})-pent-2-eno \chemfig{C(-[:150]H_3C)(-[:210]H)=C(-[:30]CH_2-CH_3)(-[:-30]H)}
e (\omterm{def:g12:stereochemistry:z-e}{E})-pent-2-eno \chemfig{C(-[:150]H_3C)(-[:210]H)=C(-[:30]H)(-[:-30]CH_2-CH_3)}.
\end{solution}

\begin{solution}{exo:g12:stereochemistry:12}
Como para o ácido tartárico, com \ce{Br} no lugar de \ce{OH} e \ce{CH3}
no lugar de \ce{COOH}: os dois \ce{Br} em cunha cheia, os dois em cunha
tracejada (um par de \omterm{def:g12:stereochemistry:enantiomer}{enantiômeros}), e um em cunha cheia e outro em cunha
tracejada, que tem um plano de simetria: a forma meso. Três
\omterm{def:g12:stereochemistry:stereoisomer}{estereoisômeros}.
\end{solution}

\begin{solution}{exo:g12:stereochemistry:13}
Metade. O químico pode separar os dois \omterm{def:g12:stereochemistry:enantiomer}{enantiômeros} (como Pasteur, por
outros meios), ou fabricar só o \omterm{def:g12:stereochemistry:enantiomer}{enantiômero} útil, com uma reação que
seja ela mesma \omterm{def:g12:stereochemistry:chiral}{quiral} (um catalisador \omterm{def:g12:stereochemistry:chiral}{quiral}, uma enzima, um material de
partida \omterm{def:g12:stereochemistry:chiral}{quiral}).
\end{solution}

\begin{solution}{exo:g12:stereochemistry:14}
Três: (2E,4E), (2Z,4Z) e (2E,4Z), sendo o último a mesma \omterm{def:g7:atoms-and-molecules:molecule}{molécula} que
(2Z,4E) lida a partir da outra ponta.
\end{solution}

\begin{solution}{exo:g12:stereochemistry:15}
$2^3 = 8$ \omterm{def:g12:stereochemistry:stereoisomer}{estereoisômeros}, isto é, 4 pares de \omterm{def:g12:stereochemistry:enantiomer}{enantiômeros}. Cada
\omterm{def:g12:stereochemistry:stereoisomer}{estereoisômero} tem um \omterm{def:g12:stereochemistry:enantiomer}{enantiômero} e $8 - 2 = 6$ \omterm{def:g12:stereochemistry:diastereomer}{diastereoisômeros}.
\end{solution}

\begin{solution}{pb:g12:stereochemistry:1}
\textbf{1.} \ce{C4H6O6}; $M = 4 \times 12.0 + 6 \times 1.0 + 6 \times 16.0
= \qty{150.0}{g/mol}$.

\textbf{2.} Dois grupos carboxila (ácido) e dois grupos hidroxila
(\omterm{def:g11:functional-groups:alcohol}{álcool}).

\textbf{3.} Os dois carbonos centrais, cada um ligado a \ce{H},
\ce{OH}, \ce{COOH} e \ce{CH(OH)COOH}.

\textbf{4.} $2^2 = 4$.

\textbf{5.} Os dois com cunhas cheias.

\textbf{6.} Os dois \ce{OH} com cunhas tracejadas; ela não pode ser
sobreposta ao L: é o ácido D-tartárico, o seu \omterm{def:g12:stereochemistry:enantiomer}{enantiômero}.

\textbf{7.} Girada em torno da ligação central de modo que os dois
\ce{OH} fiquem do mesmo lado, a \omterm{def:g7:atoms-and-molecules:molecule}{molécula} tem um plano entre os seus dois
carbonos centrais que reflete uma metade sobre a outra.

\textbf{8.} Não: é o ácido meso-tartárico.

\textbf{9.} Das quatro combinações (cheia, cheia), (tracejada,
tracejada), (cheia, tracejada), (tracejada, cheia), as duas últimas são
a mesma \omterm{def:g7:atoms-and-molecules:molecule}{molécula}, a forma meso, porque as duas metades da \omterm{def:g7:atoms-and-molecules:molecule}{molécula} são
iguais.

\textbf{10.} L e D: \omterm{def:g12:stereochemistry:enantiomer}{enantiômeros}. L e meso: \omterm{def:g12:stereochemistry:diastereomer}{diastereoisômeros}.

\textbf{11.} Também a \qty{169}{\celsius}: os \omterm{def:g12:stereochemistry:enantiomer}{enantiômeros} têm as mesmas
propriedades físicas.

\textbf{12.} Não: ele é um \omterm{def:g12:stereochemistry:diastereomer}{diastereoisômero} do L, um composto diferente
com as suas próprias propriedades.

\textbf{13.} Os receptores do nariz são \omterm{def:g12:stereochemistry:chiral}{quirais} e distinguem as imagens
no espelho; um termômetro, que não é \omterm{def:g12:stereochemistry:chiral}{quiral}, não consegue.

\textbf{14.} Não, é uma \omterm{def:g6:pure-substances-and-mixtures:mixture}{mistura} de dois compostos; como um todo, ela não
tem atividade óptica, já que os dois \omterm{def:g12:stereochemistry:enantiomer}{enantiômeros} se anulam.

\textbf{15.} Um só \omterm{def:g12:stereochemistry:enantiomer}{enantiômero}: cada cristal era feito só de L ou só de
D.

\textbf{16.} \qty{5.0}{g} de cada.

\textbf{17.} A forma meso é um composto diferente, que não faz parte da
\omterm{def:g12:stereochemistry:enantiomer}{mistura racêmica} de L e D.

\textbf{18.} Dissolvendo cada monte: as soluções desviam a luz
polarizada do mesmo tanto em sentidos opostos (e os cristais são imagens
no espelho uns dos outros).

\textbf{19.} Quatro: os seus dois \omterm{def:g12:stereochemistry:chiral}{carbonos assimétricos} carregam grupos
diferentes (\ce{Cl} num, \ce{OH} no outro), logo nenhuma combinação tem
plano de simetria, e não há duas combinações que sejam a mesma \omterm{def:g7:atoms-and-molecules:molecule}{molécula}.

\textbf{20.} Três: L, D e meso.
\end{solution}
